\documentclass[11pt,a4paper]{article}

\usepackage[T1]{fontenc}
\usepackage[utf8]{inputenc}
\usepackage[ngerman]{babel}
\usepackage{lmodern}
\usepackage{microtype}
\usepackage[a4paper,margin=2.5cm]{geometry}

\usepackage{amsmath,amssymb,amsthm,mathtools}
\usepackage{booktabs}
\usepackage{enumitem}
\usepackage[hidelinks,unicode]{hyperref}

\hypersetup{
  pdftitle={Arithmetische Präzession und Informationsgehalt},
  pdfauthor={Thomas Hoffbauer}
}

\theoremstyle{definition}
\newtheorem{definition}{Definition}[section]
\newtheorem{bemerkung}[definition]{Bemerkung}
\theoremstyle{remark}
\newtheorem{remark}[definition]{Bemerkung}

\newcommand{\N}{\mathbb{N}}
\newcommand{\R}{\mathbb{R}}

\title{Informationsgehalt, Symmetriebruch und arithmetische Präzession\\
  \large Modellierung (kein Beweisanspruch)}
\author{Thomas Hoffbauer}
\date{17.\ Juni 2026}

\begin{document}
\maketitle

\begin{abstract}
Dieses Notizblatt formalisiert eine \emph{modellhafte} Kette
Informationsgehalt $\to$ Symmetriebruch $\to$ arithmetische Perihelverschiebung
$\Pi$ auf Primvierlingsfenstern.
Es trennt drei Ebenen (Mathematik / Modellierung / Interpretation), verweist auf
den bestehenden Span-Zeugen $\mathrm{span}=10+12g$
(\texttt{Beweis\_Spanne\_EABC.tex}, \texttt{prognoseEABC.py}) und den
Projektionsrahmen (\texttt{Projektionszeuge\_Primvierling.tex},
\texttt{witness.py}).
\textbf{Kein Collatz-Anspruch}.
Der numerische Test (\texttt{collatz\_praezession\_test.py}) liefert
\emph{keine positive Evidenz} für die Kopplung $I(Q)\to\Pi(Q)$.
\end{abstract}

%======================================================================
\section{Drei Ebenen}
%======================================================================

\begin{description}[style=nextline,leftmargin=2.4cm]
  \item[\textbf{Ebene~A (fest)}]
  Markierte Startpunkte $p\equiv 5\Rightarrow\mathrm{ABCE}$,
  $p\equiv 11\Rightarrow\mathrm{CEAB}$; Primvierlingsellipse
  $(a,b,e)=(4,2,\sqrt{3}/2)$; ideale vierphasige Uhr mit $\Pi\equiv 0$
  (\texttt{Projektionszeuge\_Primvierling.tex}, \texttt{witness.py}).
  \item[\textbf{Ebene~B (Erweiterung)}]
  Integrationsstrom-Vierlinge mit Signatur \texttt{abce}/\texttt{ceab}
  und Spannenregel $\mathrm{span}(Q)=p_4-p_1=10+12g$, $g\geq 0$.
  Der Parameter $g$ zählt zusammengesetzte E-ABC-Zwischenslots auf den
  mod-$12$-Leitern (\texttt{wolfram.py}, \texttt{prognoseEABC.py}).
  \item[\textbf{Ebene~C (Modell)}]
  Informationsfunktion $I(Q)$, Kopplung $K(Q)$ und messbare Proxy-Präzession
  $\Pi(Q)$; Klein-Lesart des Symmetriebruchs (Remark unten).
\end{description}

%======================================================================
\section{Modelldefinitionen (wohldefiniert, nicht bewiesen)}
%======================================================================

Sei $Q=(p_1,p_2,p_3,p_4)$ ein Primvierlingsfenster im Integrationsstrom
(Signatur \texttt{abce} oder \texttt{ceab}) oder ein markierter Primvierling
$Q(p)=(p,p+2,p+6,p+8)$ mit $p\equiv 5,11\pmod{12}$.
Definiere
\[
  g(Q):=\frac{\mathrm{span}(Q)-10}{12}\in\N_0,
  \qquad \mathrm{span}(Q)=p_4-p_1,
\]
sofern die Spannenregel gilt (\texttt{Beweis\_Spanne\_EABC.tex}).

\begin{definition}[Informationsgehalt]
  \[
  I(Q):=\log\bigl(1+g(Q)\bigr).
  \]
  Alternativ dokumentiert: $I(Q)=\log\mathrm{span}(Q)$ oder
  $I(Q)=-\log P(Q)$ für eine Wahrscheinlichkeitsgewichtung $P$ ---
  die Wahl $\log(1+g)$ betont die diskrete Leiter-Information.
\end{definition}

\begin{definition}[Symmetrie-Kopplung]
  Mit reeller Konstante $\alpha>0$:
  \[
  K(Q):=\alpha\,I(Q)
  \qquad\text{oder}\qquad
  K(Q):=\alpha\log\bigl(1+I(Q)\bigr).
  \]
\end{definition}

\begin{definition}[Arithmetische Präzession (Proxy)]
  Für die ideale EABC-Uhr $\theta_m=\tfrac{\pi}{2}m$ gilt
  $\Pi_{\mathrm{ideal}}=0$ (\texttt{PerihelDrift.lean}).
  Als \emph{messbarer Proxy} auf Daten dient
  \[
  \Pi(Q):=\Theta_{m+4}-\Theta_m-2\pi
  \]
  unter einer dokumentierten Drift $\varepsilon(p)$, z.\,B.\
  $\varepsilon(p)=1/\log p$ (\texttt{witness.py}), oder der Betrag des
  Projektionsdefekts $|S|$ mit $S=x+z$ aus dem Heeger-Projektionszeugen.
\end{definition}

\begin{definition}[Modellgleichung]
  Mit $\beta>0$ postuliert das Gedankenmodell
  \[
  \Pi(Q)=\beta\,K(Q),
  \qquad
  \theta_{n+4}=\theta_n+2\pi+\Pi(Q).
  \]
  Dies ist \textbf{keine} Behauptung relativistischer Krümmung, sondern eine
  arithmetische Phasenkorrektur der vierphasigen Kepler-Uhr
  (\texttt{Projektionszeuge\_Primvierling.tex}, Abschnitt~Perihelverschiebung).
\end{definition}

\begin{bemerkung}[Klein-Symmetriebruch]
  Sei $S(Q)$ die Symmetriegruppe des markierten Vierlingsfensters
  (Rotationen der EABC-Tetraederdarstellung, Vertauschungen kompatibler Slots).
  In Klein-Lesart:
  \[
  I(Q)=\log\frac{|S_{\max}|}{|S(Q)|},
  \]
  wobei $S_{\max}$ die maximale symmetrische Konfiguration (ideale Ellipse,
  $\Pi=0$) bezeichnet.
  Der Symmetriebruch $I(Q)>0$ \emph{motiviert} eine nichttriviale Krümmung
  $\Pi(Q)$; die quantitative Kopplung $\Pi=\beta K$ bleibt offen.
\end{bemerkung}

%======================================================================
\section{Numerischer Test und Nichtbehauptungen}
%======================================================================

\begin{remark}[Nichtbehauptung]
  \begin{itemize}[nosep]
    \item Dieses Modell beweist weder die Collatz-Vermutung noch
    $\Einf=\emptyset$.
    \item Die Korrelation $\mathrm{corr}\bigl(I(Q),\Pi(Q)\bigr)$ auf
    Integrationsstrom-Vierlingen bis $10^6$--$10^7$ ist ein
    \emph{explorativer} Datentest (\texttt{collatz\_praezession\_test.py}),
    kein Primzahltheorem und kein Lean-Theorem.
    \item Markierte Primvierlinge $Q(p)=(p,p+2,p+6,p+8)$ haben stets
    $\mathrm{span}=8$ und $g=0$ in der Leiter-Zählung; der Span-Zeuge
    $\mathrm{span}=10+12g$ bezieht sich auf Integrationsstrom-Fenster
    (\texttt{Beweis\_Spanne\_EABC.tex}).
  \end{itemize}
\end{remark}

\begin{bemerkung}[Bestehender numerischer Zeuge]
  Die Identität $\mathrm{span}=10+12g$ ist für Integrationsstrom-Vierlinge
  bewiesen und in \texttt{wolfram.py}/\texttt{prognoseEABC.py} verifiziert
  (\texttt{collatz\_kepler\_gedankenexperiment.tex}, Tabelle numerischer Zeugen).
  Dieser Zeuge stützt die $g$-Struktur von $I(Q)$, nicht die Kopplung an $\Pi$.
\end{bemerkung}

\paragraph{Ergebnis.}
Der numerische Test liefert auf den untersuchten Integrationsstrom-Vierlingen keine positive Evidenz für eine Kopplung $I(Q)\to\Pi(Q)$.
Die gemessenen Korrelationen bleiben schwach negativ, $r\approx -0{,}18$ ($N=10^6$), $r\approx -0{,}14$ ($N=10^7$).
Damit wird die Präzessionshypothese in dieser einfachen Form nicht gestützt.

\paragraph{Einordnung.}
Das Negativergebnis betrifft nur die Kopplung $I(Q)\leftrightarrow\Pi(Q)$.
Die Span-Regel $\mathrm{span}=10+12g$ und die Ptolemaeus-/ABCE-CEAB-Geometrie bleiben unberührt.

\begin{table}[h]
  \centering
  \caption{Epistemische Taxonomie (Auszug, EABC-Projektion)}
  \begin{tabular}{@{}ll@{}}
    \toprule
    Objekt & Status \\
    \midrule
    ABCE/CEAB & Lean-bewiesen \\
    Schwerpunkt $M=p+4$ & Lean/Python/TeX \\
    Normalform $(-4,-2,2,4)$ & Lean/Python/TeX \\
    Span $10+12g$ & numerischer Zeuge \\
    Ptolemaeus / ideale Ellipse & geometrische Darstellung \\
    $\Phi_{\mathrm{pref}}$ Länge~4 & trivial \\
    $I(Q)\to\Pi(Q)$ & negativ getestet \\
    Bernoulli-Uhr & definitorische Geometrie \\
    \bottomrule
  \end{tabular}
\end{table}

\medskip
\noindent
Verwandte Texte:
\texttt{collatz\_kepler\_gedankenexperiment.tex},
\texttt{Projektionszeuge\_Primvierling.tex},
\texttt{collatz\_offene\_punkte.md} (§4 heuristisch/offen).

\end{document}
